\uuid{W3eR}
\titre{Solution particulière/solution générale}
\auteur{hanani ; mzari}
\organisation{exo7}
\licence{CC BY-NC-SA 4.0}
\datecreate{2020-01-20}
\niveau{L1}
\module{Analyse}
\chapitre{Equation différentielle}
\sousChapitre{Résolution d'équation différentielle du deuxième ordre}
\difficulte{3}
\theme{Equations du second ordre, solution particulière, solution générale}
\exo7qcmid{lille-2.equations-differentielles.029}

\contenu{
\begin{question}
\qclassification{225.02}
\qtags{motcle=solution particulière/solution générale}
On considère les équations différentielles :
$$(E_1)\; :\; y''-3y'+2y=\mathrm{e}^{2x}\quad \mbox{et}\quad (E_2)\; :\; y''-3y'+2y=2x\mathrm{e}^{x}.$$
Parmi les affirmations suivantes, cocher celles qui sont vraies :
\begin{answers}
\bad{$(E_1)$ admet une solution sous la forme $\displaystyle y_0=a\mathrm{e}^{2x}$ où $a\in \Rr$.}
\bad{$(E_2)$ admet une solution sous la forme $\displaystyle y_0=(ax+b)\mathrm{e}^{x}$ où $a,b\in \Rr$.}
\good{La solution générale de $(E_1)$ est $\displaystyle y=a\mathrm{e}^{x}+(b+x)\mathrm{e}^{2x}$, $a,b\in \Rr$.}
\good{La solution générale de $(E_2)$ est $\displaystyle y=\left(a-2x-x^2\right)\mathrm{e}^{x}+b\mathrm{e}^{2x}$, $a,b\in \Rr$.}
\end{answers}
\begin{explanations}
Les racines de l'équation caractéristique sont $1$ et $2$. Donc $(E_1)$ admet une solution particulière sous la forme $\displaystyle y_0=ax\mathrm{e}^{2x}$, car $2$ est une racine simple de l'équation caractéristique, et $(E_2)$ admet une solution particulière sous la forme $\displaystyle y_0=x(a+bx)\mathrm{e}^{x}$ car $1$ est une racine simple de l'équation caractéristique.
\end{explanations}
\end{question}
}
